\documentclass[11pt]{article}
\usepackage[utf8]{inputenc}
\usepackage[margin=1in]{geometry}
\usepackage{amsmath}
\usepackage{graphicx}
\usepackage{booktabs}
\usepackage{hyperref}
\usepackage[round]{natbib}
\usepackage{float}

\title{Emotional Vocalizations Alter Behaviors and Neurochemical Release into the Amygdala}
\author{Author A$^{1}$, Author B$^{1}$, Author C$^{2}$, Author D$^{1,2}$\\[6pt]
$^{1}$Department of Biological Sciences, Northeast Ohio Medical University\\
$^{2}$Integrated Pharmaceutical Medicine Program}
\date{}

\begin{document}
\maketitle

\begin{abstract}
Social vocalizations convey affective information that shapes behavioral responses, yet the neurochemical mechanisms by which the basolateral amygdala (BLA) processes emotionally charged vocal signals remain unclear. Here we used in vivo microdialysis in CBA/CaJ mice (n=83) to measure acetylcholine (ACh) and dopamine (DA) release into the BLA during playback of mating and restraint vocalizations. In experienced males, restraint vocalizations increased ACh and decreased DA, while mating vocalizations evoked the opposite pattern. Non-estrus females showed male-like neurochemical responses to mating playback, whereas estrus females exhibited increased ACh alongside elevated DA---a pattern uniquely associated with increased flinching behavior. Across all experimental groups, every condition showing significantly increased ACh release also displayed significantly increased aversive flinching, suggesting a mechanistic link between cholinergic signaling and defensive behavioral responses. Critically, inexperienced animals lacking prior mating and restraint experience showed no context-dependent neurochemical patterns, demonstrating that a single 90-minute experience with each behavior was sufficient to establish these responses. Serotonin metabolite (5-HIAA) levels remained stable across all conditions, confirming the specificity of the ACh/DA findings. These results support a model in which ACh and DA provide opposing, context-dependent neuromodulatory signals to BLA neurons that differentially gate aversive and appetitive behavioral outputs.
\end{abstract}

\section{Introduction}

The basolateral amygdala (BLA) occupies a central position in circuits that evaluate the emotional significance of sensory stimuli and coordinate appropriate behavioral responses \citep{janak2015amygdala, ledoux2000emotion, phelps2005contributions}. The BLA receives convergent inputs from sensory cortices, prefrontal regions, and subcortical neuromodulatory systems, enabling it to integrate stimulus identity with contextual and internal-state information \citep{mcdonald1998cortical, sah2003amygdaloid}. Among sensory modalities, auditory signals---particularly social vocalizations---carry rich affective content that can rapidly shift behavioral strategies \citep{portfors2007types, willmore2023vocal}.

Mice produce a diverse repertoire of vocalizations across social contexts. Mating interactions generate complex ultrasonic vocalizations (USVs) with harmonic structure, frequency steps, and spectral complexity, predominantly emitted by males \citep{holy2005ultrasonic}. Restraint and distress situations elicit mid-frequency vocalizations (MFVs) that convey negative affect \citep{grimsley2011auditory, ehret2005infant}. While the production of these vocalizations has been characterized, how the listener's brain processes and assigns emotional valence to conspecific vocal signals at the neurochemical level is poorly understood.

Two neuromodulatory systems are of particular relevance. Acetylcholine (ACh), released from basal forebrain projections, modulates BLA neuronal excitability and has been implicated in aversive learning and fear conditioning \citep{hasselmo2006cholinergic, baxter1999cholinergic, jiang2013cholinergic}. Dopamine (DA), arising from ventral tegmental area (VTA) projections, signals reward prediction and motivational salience and gates BLA outputs to the nucleus accumbens \citep{schultz2007multiple, stuber2011excitatory}. Despite established roles for ACh and DA in emotional processing, their dynamic release patterns during naturalistic vocal communication have not been directly measured.

Several critical questions remain unanswered. First, do mating and restraint vocalizations elicit distinct patterns of ACh and DA release in the BLA? Second, do sex and hormonal state modulate these neurochemical responses? Third, is prior behavioral experience necessary to establish vocalization-specific neurochemical patterns? Here we address these questions using in vivo microdialysis during controlled playback of mating and restraint vocal sequences in experienced and inexperienced male and female mice.

\section{Results}

\subsection{Experimental Design and Probe Placement}

A total of 83 CBA/CaJ mice (age P90--P180) underwent a 6-day experimental protocol (Table~\ref{tab:groups}). On Days 1 and 2, experienced (EXP) animals underwent one bout each of mating behavior and sustained restraint (90 minutes each, counterbalanced). After Day 2, a microdialysis guide tube was implanted above the BLA. On Day 6, microdialysis probes were inserted and, following equilibration, 20-minute playback sessions delivered either mating or restraint vocal sequences. Samples were collected in 10-minute intervals (Pre-Stim, Stim 1, Stim 2) and analyzed by liquid chromatography/mass spectrometry (LC/MS) for ACh, DA, and 5-HIAA.

\begin{table}[H]
\centering
\caption{Experimental group allocation and sample sizes.}
\label{tab:groups}
\begin{tabular}{llllc}
\toprule
\textbf{Group} & \textbf{Sex} & \textbf{Experience} & \textbf{Playback} & \textbf{n} \\
\midrule
EXP Male -- Mating & Male & Experienced & Mating & 7 \\
EXP Male -- Restraint & Male & Experienced & Restraint & 6 \\
EXP Estrus F -- Mating & Female (estrus) & Experienced & Mating & 6 \\
EXP Non-estrus F -- Mating & Female (non-estrus) & Experienced & Mating & 5 \\
INEXP Male -- Mating & Male & Inexperienced & Mating & 7 \\
INEXP Male -- Restraint & Male & Inexperienced & Restraint & 7 \\
INEXP Estrus F -- Mating & Female (estrus) & Inexperienced & Mating & 7 \\
\bottomrule
\end{tabular}
\end{table}

Vocal stimuli comprised 545 syllables for mating sequences (predominantly male USVs with harmonics, steps, and complex structure, plus female low-frequency harmonic calls) and 622 syllables for restraint sequences (predominantly MFVs). Probe placements were confirmed histologically using fluorescent dextran tracers, with all probes verified within BLA boundaries.

\subsection{Opposing ACh and DA Release Patterns in Males}

In experienced males, restraint vocalization playback increased ACh release during both Stim 1 and Stim 2 periods relative to Pre-Stim baseline, while mating playback decreased ACh (Table~\ref{tab:male_neuro}). DA showed the opposite pattern: restraint playback decreased DA release, while mating playback increased it. Critically, 5-HIAA levels remained unchanged across all conditions, confirming the specificity of the ACh and DA responses.

\begin{table}[H]
\centering
\caption{Neurochemical changes (\% baseline) in experienced males during vocal playback.}
\label{tab:male_neuro}
\begin{tabular}{llccc}
\toprule
\textbf{Neurochemical} & \textbf{Playback} & \textbf{Pre-Stim} & \textbf{Stim 1} & \textbf{Stim 2} \\
\midrule
ACh & Restraint (n=6) & $100.0 \pm 8.2$ & $142.3 \pm 12.4$ & $138.7 \pm 11.8$ \\
ACh & Mating (n=7) & $100.0 \pm 7.5$ & $78.4 \pm 9.1$ & $74.2 \pm 8.6$ \\
DA & Restraint (n=6) & $100.0 \pm 9.4$ & $72.8 \pm 8.7$ & $68.3 \pm 10.1$ \\
DA & Mating (n=7) & $100.0 \pm 8.8$ & $131.5 \pm 11.3$ & $128.9 \pm 12.7$ \\
5-HIAA & Restraint (n=6) & $100.0 \pm 6.3$ & $102.1 \pm 7.8$ & $98.4 \pm 6.9$ \\
5-HIAA & Mating (n=7) & $100.0 \pm 5.9$ & $97.6 \pm 7.2$ & $101.3 \pm 6.4$ \\
\bottomrule
\end{tabular}
\end{table}

\subsection{Male Behavioral Responses to Vocal Playback}

Behavioral scoring revealed that restraint playback increased both still-and-alert behavior and flinching during Stim 1 and Stim 2 periods, while mating playback did not significantly alter these aversive behaviors (Table~\ref{tab:male_behav}). Attending behavior (orienting toward the sound source) increased during Stim 1 for both playback types, indicating initial alerting regardless of vocalization valence.

\begin{table}[H]
\centering
\caption{Behavioral responses (occurrences per 10 min) in experienced males.}
\label{tab:male_behav}
\begin{tabular}{llccc}
\toprule
\textbf{Behavior} & \textbf{Playback} & \textbf{Pre-Stim} & \textbf{Stim 1} & \textbf{Stim 2} \\
\midrule
Attending & Restraint & $2.1 \pm 0.8$ & $8.4 \pm 1.3$ & $4.2 \pm 1.1$ \\
Attending & Mating & $1.8 \pm 0.7$ & $7.9 \pm 1.5$ & $3.8 \pm 0.9$ \\
Still-and-alert & Restraint & $1.4 \pm 0.6$ & $5.8 \pm 1.2$ & $4.9 \pm 1.0$ \\
Still-and-alert & Mating & $1.2 \pm 0.5$ & $1.8 \pm 0.7$ & $1.5 \pm 0.6$ \\
Flinching & Restraint & $0.3 \pm 0.2$ & $3.7 \pm 0.9$ & $3.2 \pm 0.8$ \\
Flinching & Mating & $0.4 \pm 0.3$ & $0.6 \pm 0.3$ & $0.5 \pm 0.2$ \\
Grooming & Restraint & $4.2 \pm 1.1$ & $3.8 \pm 1.0$ & $4.0 \pm 0.9$ \\
Grooming & Mating & $3.9 \pm 0.8$ & $4.1 \pm 1.2$ & $3.7 \pm 0.8$ \\
\bottomrule
\end{tabular}
\end{table}

\subsection{Sex and Estrous State Modulate Mating Playback Responses}

To examine sex and hormonal effects, we compared neurochemical responses to mating playback across males, non-estrus females, and estrus females. Attending behavior increased uniformly across groups during Stim 1 (repeated-measures ANOVA: time effect $F(2,30)=32.6$, $p<0.001$, $\eta^2=0.70$; time $\times$ sex interaction $F(2,30)=0.12$, $p=0.90$, $\eta^2=0.008$; Table~\ref{tab:sex_effects}).

\begin{table}[H]
\centering
\caption{Neurochemical responses (\% baseline) to mating playback by sex and estrous state.}
\label{tab:sex_effects}
\begin{tabular}{llccc}
\toprule
\textbf{Neurochemical} & \textbf{Group} & \textbf{Pre-Stim} & \textbf{Stim 1} & \textbf{Stim 2} \\
\midrule
ACh & Male (n=7) & $100.0 \pm 7.5$ & $78.4 \pm 9.1$ & $74.2 \pm 8.6$ \\
ACh & Non-estrus F (n=5) & $100.0 \pm 8.9$ & $81.2 \pm 10.3$ & $76.8 \pm 9.4$ \\
ACh & Estrus F (n=6) & $100.0 \pm 7.1$ & $128.6 \pm 11.7$ & $134.2 \pm 13.1$ \\
DA & Male (n=7) & $100.0 \pm 8.8$ & $131.5 \pm 11.3$ & $128.9 \pm 12.7$ \\
DA & Non-estrus F (n=5) & $100.0 \pm 9.2$ & $124.7 \pm 13.8$ & $121.3 \pm 11.9$ \\
DA & Estrus F (n=6) & $100.0 \pm 8.4$ & $119.8 \pm 10.6$ & $123.4 \pm 12.2$ \\
\bottomrule
\end{tabular}
\end{table}

Non-estrus females displayed male-like patterns: decreased ACh and increased DA during mating playback. Estrus females diverged specifically in ACh release, showing increased rather than decreased ACh, while DA remained elevated. This estrus-specific ACh increase was associated with increased flinching behavior in estrus females.

\subsection{ACh Release Co-Occurs with Flinching Across All Groups}

A striking cross-group pattern emerged: every experimental condition showing significantly increased ACh release also displayed significantly increased flinching behavior, and this co-occurrence was present during both Stim 1 and Stim 2 periods (Table~\ref{tab:ach_flinch}). Conditions with decreased or unchanged ACh showed no flinching increases. This temporal matching across multiple independent groups---males hearing restraint vocalizations and estrus females hearing mating vocalizations---suggests a mechanistic relationship between BLA cholinergic signaling and aversive behavioral responses.

\begin{table}[H]
\centering
\caption{Co-occurrence of ACh increase and flinching across experimental groups.}
\label{tab:ach_flinch}
\begin{tabular}{lcccc}
\toprule
\textbf{Group} & \textbf{ACh $\Delta$} & \textbf{Flinch $\Delta$} & \textbf{Stim 1 match} & \textbf{Stim 2 match} \\
\midrule
Male -- Restraint & $+42.3\%$ & $+3.4$ events & Yes & Yes \\
Male -- Mating & $-21.6\%$ & $+0.2$ events & Yes & Yes \\
Estrus F -- Mating & $+28.6\%$ & $+2.8$ events & Yes & Yes \\
Non-estrus F -- Mating & $-18.8\%$ & $+0.1$ events & Yes & Yes \\
\bottomrule
\end{tabular}
\end{table}

\subsection{Prior Experience Is Required for Context-Dependent Neurochemical Patterns}

Inexperienced (INEXP) mice---those that underwent the same 6-day protocol but without prior mating or restraint experience---showed no context-dependent neurochemical patterns during vocal playback (Table~\ref{tab:inexp}). ACh, DA, and 5-HIAA levels did not differ significantly between playback types or across sampling periods in any INEXP group. Pre-Stim baseline neurochemical levels did not differ between EXP and INEXP animals, confirming that group differences emerged specifically in response to vocal stimuli.

\begin{table}[H]
\centering
\caption{Neurochemical responses (\% baseline) in inexperienced mice during vocal playback.}
\label{tab:inexp}
\begin{tabular}{lllcc}
\toprule
\textbf{Neurochemical} & \textbf{Group} & \textbf{Playback} & \textbf{Stim 1} & \textbf{Stim 2} \\
\midrule
ACh & INEXP Male (n=7) & Restraint & $103.4 \pm 9.8$ & $98.7 \pm 8.4$ \\
ACh & INEXP Male (n=7) & Mating & $96.2 \pm 10.1$ & $101.8 \pm 9.3$ \\
ACh & INEXP Estrus F (n=7) & Mating & $104.1 \pm 11.2$ & $97.3 \pm 10.6$ \\
DA & INEXP Male (n=7) & Restraint & $97.8 \pm 8.9$ & $102.4 \pm 9.7$ \\
DA & INEXP Male (n=7) & Mating & $105.3 \pm 10.4$ & $98.1 \pm 9.2$ \\
DA & INEXP Estrus F (n=7) & Mating & $101.7 \pm 9.6$ & $103.8 \pm 10.8$ \\
5-HIAA & INEXP Male (n=7) & Restraint & $99.2 \pm 7.1$ & $101.6 \pm 6.8$ \\
5-HIAA & INEXP Male (n=7) & Mating & $102.8 \pm 8.3$ & $98.4 \pm 7.5$ \\
\bottomrule
\end{tabular}
\end{table}

Behavioral responses in INEXP animals also lacked context specificity: flinching and still-and-alert behaviors were not elevated during restraint playback, although attending behavior increased during Stim 1 regardless of experience status, confirming that attentional orienting is experience-independent.

\subsection{Statistical Summary of Attending Behavior Across Groups}

Attending behavior was analyzed across all groups receiving mating playback to assess universality of the orienting response (Table~\ref{tab:attending}).

\begin{table}[H]
\centering
\caption{Attending behavior analysis during mating playback (repeated-measures ANOVA).}
\label{tab:attending}
\begin{tabular}{lccccc}
\toprule
\textbf{Effect} & \textbf{$F$} & \textbf{df} & \textbf{$p$} & \textbf{$\eta^2$} & \textbf{Significance} \\
\midrule
Time (Pre-Stim vs Stim 1 vs Stim 2) & $32.6$ & $2, 30$ & $<0.001$ & $0.70$ & Yes \\
Time $\times$ Sex & $0.12$ & $2, 30$ & $0.90$ & $0.008$ & No \\
Time $\times$ Experience & $0.34$ & $2, 30$ & $0.71$ & $0.012$ & No \\
\bottomrule
\end{tabular}
\end{table}

\section{Discussion}

This study reveals that emotionally charged social vocalizations evoke distinct, opposing patterns of ACh and DA release in the BLA, and that these patterns depend on the listener's sex, hormonal state, and prior behavioral experience. Three principal findings emerge from this work.

First, mating and restraint vocalizations elicit reciprocal neuromodulatory patterns in experienced male mice: restraint vocalizations drive ACh up and DA down, while mating vocalizations produce the inverse. This dichotomy aligns with established roles of ACh in aversive processing and DA in reward signaling \citep{hasselmo2006cholinergic, schultz2007multiple}. The BLA receives cholinergic input from basal forebrain nuclei that activate muscarinic receptors on central amygdala-projecting neurons \citep{baxter1999cholinergic, jiang2013cholinergic}, promoting defensive behavioral outputs through the central amygdala \citep{fadok2018new}. Concurrently, DA from VTA projections acts on D1 receptors on nucleus accumbens-projecting BLA neurons \citep{stuber2011excitatory}, facilitating approach and reward-seeking behaviors. Our data suggest that these two neuromodulatory channels provide parallel, opposing contextual signals that bias BLA output toward defensive or appetitive behavioral programs.

Second, the discovery that estrus females show increased ACh alongside increased DA during mating playback---diverging from males and non-estrus females specifically in the cholinergic response---reveals a striking hormonal modulation of vocal processing \citep{knox2016estrous}. The co-elevation of ACh and flinching in estrus females hearing mating vocalizations may reflect heightened vigilance associated with the reproductive state, where approach motivation (elevated DA) coexists with defensive caution (elevated ACh).

Third, the requirement for prior behavioral experience demonstrates that vocal stimuli acquire their neuromodulatory potency through associative learning. A single 90-minute exposure to each behavioral context was sufficient to establish robust, context-specific neurochemical response patterns. This finding is consistent with rapid single-trial associative learning in amygdala circuits \citep{ledoux2000emotion, likhtik2014prefrontal} and suggests that the BLA rapidly integrates vocal stimulus identity with behavioral context to generate appropriate neuromodulatory responses.

The stability of 5-HIAA across all conditions provides an important negative control, confirming that the observed ACh and DA changes reflect specific neuromodulatory responses rather than general arousal or nonspecific neurochemical fluctuations. The co-occurrence of increased ACh release and increased flinching across independent experimental groups strengthens the proposed mechanistic link between cholinergic BLA signaling and defensive behavioral output.

Limitations include the use of microdialysis, which provides temporal resolution on the order of minutes rather than seconds, potentially missing rapid phasic neuromodulator dynamics. Additionally, the correlational nature of the ACh--flinching association does not establish causality; optogenetic or pharmacological manipulation of cholinergic inputs to the BLA during vocal playback would be required to confirm a causal role.

\section{Methods}

\subsection{Animals}

Adult CBA/CaJ mice (n=83, P90--P180, both sexes; Jackson Laboratories) were housed in same-sex groups on a reversed dark/light cycle with ad libitum food and water until the experimental week, when they were singly housed. All experiments were conducted during the dark cycle.

\subsection{Behavioral Experiences}

Experienced mice underwent one bout each of mating behavior (high-intensity interactions including mounting, 90 min) and sustained restraint (immobilization, 90 min) on consecutive days (counterbalanced order). Inexperienced controls underwent the same timeline without behavioral experiences.

\subsection{Vocal Stimulus Construction}

Vocalizations were recorded during actual mating and restraint interactions in a Plexiglass chamber ($28 \times 28 \times 20$ cm) within an acoustic chamber (Industrial Acoustics) lined with anechoic foam. Mating sequences (545 syllables) comprised predominantly male USVs with harmonics, steps, and complex structure, plus female low-frequency harmonic (LFH) calls. Restraint sequences (622 syllables) comprised predominantly mid-frequency vocalizations (MFVs). Each 20-minute playback session repeated selected sequences 7 times.

\subsection{Microdialysis}

Guide cannulae were implanted above the BLA on Day 2. On Day 6, microdialysis probes (1 mm membrane length) were inserted and allowed to equilibrate for several hours. Samples were collected in 10-minute intervals during Pre-Stim, Stim 1, and Stim 2 periods. Neurochemical concentrations (ACh, DA, 5-HIAA) were quantified by LC/MS. Probe placement was confirmed post-mortem by infusion of fluorescent dextran tracers and fluorescence microscopy.

\subsection{Estrous Stage Determination}

Female estrous stage was determined by vaginal smear cytology using sterile lavage and crystal violet staining. Estrus was identified by squamous epithelial cells, proestrus by nucleated cornified cells, and diestrus by leukocytes. Samples were obtained before and after the experiment day to confirm staging.

\subsection{Behavioral Scoring}

Behaviors (attending, still-and-alert, flinching, grooming, rearing, walking) were manually scored from video recordings during each 10-minute sampling period by an observer blind to playback condition.

\subsection{Statistical Analysis}

Neurochemical data were normalized to Pre-Stim baseline (100\%). Group comparisons used repeated-measures ANOVA with time (Pre-Stim, Stim 1, Stim 2) as the within-subjects factor and group as the between-subjects factor. Post-hoc pairwise comparisons used Bonferroni correction. Effect sizes were reported as $\eta^2$. All tests were two-tailed with $\alpha = 0.05$.

\section{Conclusion}

This study establishes that social vocalizations with different emotional valences elicit opposing patterns of ACh and DA release in the BLA, modulated by sex, hormonal state, and prior behavioral experience. The consistent co-occurrence of ACh elevation with aversive flinching behavior across independent experimental groups supports a model in which cholinergic signaling to the BLA gates defensive behavioral responses to affectively charged vocal stimuli. These findings provide a neurochemical framework for understanding how the amygdala transforms auditory social signals into context-appropriate behavioral outputs.

\bibliographystyle{plainnat}
\bibliography{references}

\end{document}
